\documentclass[11pt]{article}
\usepackage[a4paper,margin=26mm]{geometry}
\usepackage{amsmath,amssymb,amsthm,hyperref}
\newtheorem{theorem}{Theorem}
\newtheorem{lemma}{Lemma}
\title{Every sufficiently large quadratic size is possible for one die}
\author{Research corollary; two independent full reviews completed}
\date{30 September 2026}
\begin{document}\maketitle

\begin{theorem}
There is an absolute constant $C$ such that, for every $n\ge2$ and
\emph{every integer} $K\ge Cn^2$, there is a tie-free permutation-fair
collection of $n$ equiprobable finite dice in which a specified die has
exactly $K$ faces. The other $n-1$ dice can all have the same number
of faces, which is not bounded quantitatively here.
\end{theorem}

This strengthens the reviewed existence of some exceptional size
$K=O(n^2)$. In conjunction with the reviewed universal lower bound
$K\ge cn^2$, it locates the entire possible-size spectrum up to
absolute constant factors. No divisibility condition is imposed on
$K$: denominator clearing occurs only on the other dice.

\begin{lemma}
For every $m\ge1$ and every integer $K\ge C_0(m+1)^2$, there is a
real equal-weight quadrature for the uniform probability measure on
$[0,1]$, exact through degree $m$, with exactly $K$ distinct nodes
strictly inside $(0,1)$.
\end{lemma}
\begin{proof}
Gilboa--Peled, \emph{Chebyshev-type Quadratures for Doubling Weights},
Theorem 1.1, bounds both $N_w(d)$ and $\overline N_w(d)$ by a constant
multiple of their local-mass quantity $R_w(d)$.
Their definition of $\overline N_w(d)$, on pages 2--3, requires a rule
with exactly $N$ nodes for \emph{every} $N\ge\overline N_w(d)$.
For the constant weight on $[-1,1]$, the smallest window mass is of
order $d^{-2}$; hence $\overline N_w(d)\le C_1d^2$.
Affine rescaling gives the same statement on $[0,1]$.
The primary text and theorem are available at
\url{https://arxiv.org/pdf/1507.01505}, pages 2--3.

Apply that result at degree $2m+2$, choosing the prescribed $K$.
Increase $C_0$ so that it also ensures $K\ge m^2$.
The resulting rule has at least $m+2$ distinct support points: if it
had at most $m+1$, the square of the support-vanishing polynomial
would have degree at most $2m+2$, zero quadrature sum, and strictly
positive integral. Thus at least $m$ distinct support points are
interior.

Select one node at each of $m$ distinct interior points as pivots.
The Jacobian of the first $m$ moment sums in these coordinates is
an invertible scaled Vandermonde matrix. By the implicit function
theorem, small arbitrary perturbations of the remaining coordinates
can be compensated by the pivots while preserving all these moments.
Move free endpoint coordinates slightly inward. Pivot coordinates
remain interior and distinct. Free--free collisions are proper
hyperplanes. At any initial collision of a free node and a pivot,
the matching pivot has derivative $-1$ when that free node varies,
so their difference has nonzero derivative $-2$. These collisions
are therefore proper analytic zero sets. Initially absent collisions
remain absent after taking a small neighborhood.
A finite union of proper analytic zero sets does not fill the open
set of permitted inward perturbations. A collision-free choice,
followed by sorting, gives exactly $K$ distinct interior nodes.
\end{proof}

\begin{proof}[Proof of the theorem]
Set $m=n-1$ and use the preceding lemma at the prescribed $K$.
The reviewed proof in
\path{../../size_bounds/polynomial_exceptional_die.tex} now applies
without repeating, adding, or removing any distinguished nodes.
We detail why its quantifiers permit this application.

That proof requires a degree-$m$ equal-weight real rule with
$K\ge m^2$ distinct interior nodes, and no upper bound on $K$.
It selects disjoint sets $Q_1,\ldots,Q_m$, each containing $m$ of
those nodes, and rational moment bumps with disjoint supports around
the nodes. The node locations are approximated by rationals.
For each old die, an invertible $m$-by-$m$ Vandermonde system finds
its correction coefficients. These coefficients tend to zero as
the rational approximation approaches the real rule, so all corrected
densities remain positive. The disjoint-support response identity
then makes every full order have probability $1/n!$.

The distinguished distribution is still uniform on the original
number $K$ of rational locations. All other distributions have
positive rational piecewise-constant densities. Partitioning at their
breakpoints and at the $K$ distinguished locations, then replacing
each interval by a finite fair block, preserves the complete order
law. A common denominator $D$ for the interval masses and a common
$F$-face fair block give exactly $FD$ faces for each old die.
One distinguished face is inserted at each of the $K$ selected
locations. Assigning globally distinct integer labels finishes the
construction with exactly the requested distinguished size.
\end{proof}

The finite realization may require very large $F$ and $D$. This
corollary gives neither a small total-size construction nor simultaneous
control of the sizes of several specified dice.
\end{document}
